\documentclass[conference]{IEEEtran}

\usepackage{amsmath,amssymb,amsfonts}
\usepackage{algorithmic}
\usepackage{algorithm}
\usepackage{graphicx}
\usepackage{textcomp}
\usepackage{booktabs}
\usepackage{cite}
\usepackage{hyperref}
\usepackage{xcolor}
\usepackage{amsthm}
\usepackage{CJKutf8}
\usepackage{url}
\usepackage{multirow}
\usepackage{enumitem}

\newtheorem{theorem}{定理}
\newtheorem{definition}{定義}
\newtheorem{lemma}{補題}
\newtheorem{corollary}{系}
\newtheorem{proposition}{命題}

\begin{document}
\begin{CJK}{UTF8}{ipxm}

\title{DenseInterp: 密集アイテムセットパターンの\\LLM による構造化自然言語解釈フレームワーク}

\author{
\IEEEauthorblockN{Anonymous Author(s)}
}

\maketitle

% =====================================================================
\begin{abstract}
頻出アイテムセットマイニングは大規模トランザクションデータから有用なパターンを発見する基盤技術であるが，
検出パターンの解釈はドメイン専門家に依存し，分析のボトルネックとなっている．
本論文では，密集アイテムセット（スライディングウィンドウ内で共起頻度が閾値を超えるアイテムの組み合わせ）の
検出結果を大規模言語モデル（LLM）により自然言語解釈に変換するフレームワーク DenseInterp を提案する．
提案手法は4つの主要コンポーネントから構成される：
(1) パターンの時間的構造を定量化する特徴量抽出器，
(2) 解釈タスクを単一パターン・比較・バッチの3層に分解する構造化プロンプト生成器，
(3) LLMバックエンドの抽象化レイヤ，
(4) 解釈品質を多軸で評価するメトリクス体系．
3つの実データセット（Online Retail, Dunnhumby, Chicago Crime）上で
GPT-4o, Claude 3.5 Sonnet, Llama-3-70B の3モデルを用いた実験により，
提案フレームワークが平均 0.82 の事実整合性スコアと 0.76 の仮説妥当性スコアを達成し，
非構造化プロンプトに対して事実整合性で $+$18.3\%，仮説妥当性で $+$23.7\% の改善を示すことを確認した．
さらに，ドメイン専門家3名による人間評価では，
DenseInterp の解釈が専門家自身の解釈と平均 78.4\% の一致率を示した．
\end{abstract}

% =====================================================================
\section{はじめに}
\label{sec:intro}

頻出アイテムセットマイニング~\cite{agrawal1994fast, han2000mining}は，
トランザクションデータベースから高頻度で共起するアイテムの組み合わせを発見する
データマイニングの基本タスクである．
近年，スライディングウィンドウに基づく密集区間検出手法が提案され，
パターンの時間的局在性を捉える能力が向上している~\cite{mannila1997discovery}．

しかし，検出パターンの\emph{解釈}は依然として専門家に依存している．
例えば，「アイテム $\{A, B, C\}$ がトランザクション $[50, 120]$ で密集」
という検出結果から，なぜこのパターンが生じたのか，
ビジネス的にどのような意味を持つのかを判断するには
ドメイン知識と分析経験が不可欠である．
Fournier-Viger ら~\cite{fournier2014spmf}が指摘するように，
パターンマイニングツールは大量のパターンを出力するが，
その解釈支援機能は極めて限定的である．

大規模言語モデル（LLM）~\cite{brown2020gpt3, openai2023gpt4}の発展により，
構造化データの自然言語説明の自動生成が現実的な技術となりつつある．
しかし，マイニング結果を単純にLLMに入力するだけでは，
時間的構造の無視，数値的事実との不整合（ハルシネーション），
ドメイン非依存な汎用的記述に留まる，といった問題が生じる~\cite{ji2023hallucination}．

本論文では，これらの課題を解決するフレームワーク DenseInterp を提案する．
主要な貢献は以下の4点である：

\begin{itemize}[leftmargin=*]
    \item \textbf{構造化特徴量抽出}: パターンの時間的構造（区間数・カバー率・
          ギャップ分布・周期性）を情報理論的指標を含む12次元特徴ベクトルとして定量化し，
          LLMが解釈に必要な文脈情報を体系的に提供する手法
    \item \textbf{3層プロンプト設計}: 単一パターン解釈・複数パターン比較・
          バッチ要約の3層構造により，粒度の異なる解釈ニーズに対応する
          テンプレート体系の設計と，各層の最適プロンプト戦略の同定
    \item \textbf{多軸評価メトリクス}: 事実整合性・仮説妥当性・
          時間記述精度・表現多様性の4軸で解釈品質を定量化する
          評価フレームワークの提案と，人間評価との相関分析
    \item \textbf{3モデル$\times$3データセットの実証実験}: GPT-4o, Claude 3.5 Sonnet,
          Llama-3-70B を用いた体系的な比較実験と，ドメイン専門家による人間評価
\end{itemize}

% =====================================================================
\section{関連研究}
\label{sec:related}

\subsection{頻出パターンマイニングと解釈可能性}

Apriori~\cite{agrawal1994fast}とFP-Growth~\cite{han2000mining}は
頻出アイテムセット発見の基盤アルゴリズムである．
エピソードマイニング~\cite{mannila1997discovery}はイベント列からの
時間的パターン発見を扱い，Webb~\cite{webb2007discovering}は
統計的有意性による偽陽性パターンの排除を提案した．

パターンマイニングの解釈可能性に関しては，
SPMF~\cite{fournier2014spmf}がパターン可視化機能を提供するが，
自然言語による自動解釈は未対応である．
Vreeken ら~\cite{vreeken2011krimp}のKRIMP はMDL原理による
パターン集合の圧縮要約を提案したが，個別パターンの
意味論的解釈は扱っていない．
本研究は，密集区間という時間的構造を持つパターンに特化した
自然言語解釈を提供する点で，既存のパターン要約手法と相補的である．

\subsection{説明可能AI（XAI）とデータ分析の自動解釈}

LIME~\cite{ribeiro2016lime}やSHAP~\cite{lundberg2017shap}に代表されるXAI手法は，
機械学習モデルの予測根拠を説明する技術として広く研究されている．
しかし，これらはモデルの入出力関係の局所的説明に焦点を当てており，
教師なし学習であるパターンマイニングの結果解釈には直接適用できない．

自動レポート生成の分野では，
データからの自然言語記述生成（Data-to-Text）~\cite{reiter2007architecture}が
気象データや財務データを対象に研究されてきた．
Shen ら~\cite{shen2023data}はLLMを用いたデータ分析の自動化を
調査し，構造化データの理解においてプロンプト設計が
品質に大きく影響することを示した．
本研究はこの知見を密集パターン解釈に特化して応用する．

\subsection{LLM による構造化データ理解}

GPT-3~\cite{brown2020gpt3}以降，LLMの構造化データへの応用が
活発に研究されている．
Chain-of-Thought (CoT) プロンプティング~\cite{wei2022chain}は
段階的推論を促し，数値データの解釈精度を向上させる．
Hegselmann ら~\cite{hegselmann2023tabllm}はTabLLMにおいて
表形式データのテキスト直列化がLLM理解に与える影響を分析し，
構造化シリアライゼーションの重要性を示した．

LLMのハルシネーション問題~\cite{ji2023hallucination}は
数値データの解釈において特に深刻である．
本研究では，パターンの数値的特徴量を明示的にプロンプトに含め，
CoT推論を組み合わせることで事実整合性を向上させる．

Narayan ら~\cite{narayan2018dont}の要約生成技術は
長文テキストの圧縮に有効であるが，
構造化パターンの意味論的要約には直接適用が難しい．
Zhao ら~\cite{zhao2023survey}のサーベイはLLMの包括的な能力と
限界を整理しており，本研究の設計方針に影響を与えた．

\subsection{プロンプトエンジニアリング}

プロンプト設計は LLM の出力品質に決定的な影響を与える~\cite{liu2023summary}．
Few-shot 学習~\cite{brown2020gpt3}は少数の例示により
タスク理解を促進し，CoT~\cite{wei2022chain}は
中間推論ステップの明示化により複雑な推論精度を改善する．
Self-consistency~\cite{wang2023selfconsistency}は
複数の推論パスからの多数決により安定性を向上させる．
本研究では，これらの技法を密集パターン解釈に最適化した
3層プロンプトテンプレートとして体系化する．

% =====================================================================
\section{問題定義}
\label{sec:problem}

\begin{definition}[密集アイテムセット]
\label{def:dense}
トランザクションデータベース $\mathcal{D} = \{T_0, T_1, \ldots, T_{n-1}\}$，
ウィンドウサイズ $W \in \mathbb{Z}^+$，閾値 $\theta \in \mathbb{Z}^+$ が与えられたとき，
アイテムセット $I \subseteq \mathcal{I}$ の\emph{密集区間} $[s, e]$（$0 \leq s \leq e \leq n - W$）とは，
すべての $l \in [s, e]$ について
$\text{supp}_W(I, l) = |\{T_j : j \in [l, l+W), I \subseteq T_j\}| \geq \theta$
を満たす極大連続区間である．
アイテムセット $I$ の密集区間集合を $\mathcal{S}(I) = \{[s_1, e_1], \ldots, [s_k, e_k]\}$
（$e_j < s_{j+1}$ for all $j$）と表記する．
\end{definition}

\begin{definition}[パターン特徴量ベクトル]
\label{def:features}
密集パターン $P = (I, \mathcal{S}(I))$（$\mathcal{S}(I) = \{[s_j, e_j]\}_{j=1}^{k}$）に対して，
特徴量ベクトル $\mathbf{f}(P) \in \mathbb{R}^{12}$ を以下のように定義する：

\noindent\textbf{基本特徴量}（4次元）：
\begin{itemize}[leftmargin=*,nosep]
    \item カバー率: $\rho = \sum_{j=1}^{k}(e_j - s_j + 1) / (n - W + 1)$
    \item 区間数: $k = |\mathcal{S}(I)|$
    \item アイテムセットサイズ: $|I|$
    \item 平均区間長: $\bar{L} = \frac{1}{k}\sum_{j=1}^{k}(e_j - s_j + 1)$
\end{itemize}

\noindent\textbf{ギャップ特徴量}（3次元，$k \geq 2$ の場合）：
\begin{itemize}[leftmargin=*,nosep]
    \item 最大ギャップ: $g_{\max} = \max_{j=1}^{k-1}(s_{j+1} - e_j - 1)$
    \item 平均ギャップ: $\bar{g} = \frac{1}{k-1}\sum_{j=1}^{k-1}(s_{j+1} - e_j - 1)$
    \item ギャップ変動係数: $\text{CV}_g = \sigma_g / \bar{g}$
\end{itemize}

\noindent\textbf{時間分布特徴量}（3次元）：
\begin{itemize}[leftmargin=*,nosep]
    \item 開始位置（正規化）: $\hat{s}_1 = s_1 / (n - W)$
    \item 終了位置（正規化）: $\hat{e}_k = e_k / (n - W)$
    \item 時間的集中度: $\text{TC} = 1 - (\hat{e}_k - \hat{s}_1)$ if $k \geq 1$
\end{itemize}

\noindent\textbf{分類特徴量}（2次元）：
\begin{itemize}[leftmargin=*,nosep]
    \item 時間パターン分類（one-hot の代わりに序数）:
        $\tau = \begin{cases}
            1 & k = 1 \text{ (単一連続)} \\
            2 & 2 \leq k \leq 3 \text{ (断続的)} \\
            3 & k \geq 4 \text{ (高頻度断続)}
        \end{cases}$
    \item 周期性指標: $\pi = 1 - \text{CV}_g$ ($k \geq 3$ のとき有効, それ以外は $0$)
\end{itemize}
\end{definition}

\begin{definition}[パターン解釈タスク]
\label{def:interpretation}
密集パターン $P = (I, \mathcal{S}(I))$，特徴量ベクトル $\mathbf{f}(P)$，
オプショナルなアイテムメタデータ $M: \mathcal{I} \to \text{String}$，
およびドメインコンテキスト $\mathcal{C}$（業種・時間粒度等の記述）が与えられたとき，
パターン解釈関数 $\Phi: (P, \mathbf{f}(P), M, \mathcal{C}) \to R$ は
以下の4要素を含む構造化テキスト $R$ を生成する：
\begin{enumerate}[nosep]
    \item \textbf{要約} $R_{\text{sum}}$: パターンの1--2文要約
    \item \textbf{時間的記述} $R_{\text{temp}}$: 密集区間の時間的特徴の記述
    （$\mathbf{f}(P)$ の数値と矛盾しないこと）
    \item \textbf{ビジネス仮説} $R_{\text{hyp}}$: パターンが生じた原因の仮説
    （$\mathcal{C}$ に基づくこと）
    \item \textbf{確信度} $c \in [0, 1]$: 解釈の確からしさの自己評価
\end{enumerate}
\end{definition}

\begin{definition}[解釈品質メトリクス]
\label{def:metrics}
解釈結果の集合 $\mathcal{R} = \{R_1, \ldots, R_m\}$ に対して，
以下の4軸で品質を評価する：

\noindent\textbf{事実整合性}（Factual Consistency, FC）：
$\text{FC}(R_i) = \frac{|\text{Claims}(R_i) \cap \text{Facts}(\mathbf{f}(P_i))|}{|\text{Claims}(R_i)|}$
ここで $\text{Claims}(R_i)$ は解釈中の数値的主張の集合，
$\text{Facts}(\mathbf{f}(P_i))$ は特徴量から検証可能な事実の集合．

\noindent\textbf{仮説妥当性}（Hypothesis Plausibility, HP）：
$\text{HP}(R_i) \in [0, 1]$，ドメイン専門家による評価（5段階を正規化）．

\noindent\textbf{時間記述精度}（Temporal Description Accuracy, TDA）：
$\text{TDA}(R_i) = 1 - \frac{|\hat{\rho}_i - \rho_i|}{\rho_i}$
ここで $\hat{\rho}_i$ は解釈文中で言及されたカバー率の推定値．

\noindent\textbf{表現多様性}（Expression Diversity, ED）：
$\text{ED}(\mathcal{R}) = 1 - \frac{2}{m(m-1)}\sum_{i<j}\text{BLEU}(R_i, R_j)$
同一テンプレートからの出力が画一的でないかを測定．
\end{definition}

% =====================================================================
\section{提案手法: DenseInterp}
\label{sec:method}

提案フレームワーク DenseInterp は4つのモジュールから構成される
（図~\ref{fig:architecture}）．

\begin{figure}[t]
    \centering
    \fbox{\parbox{0.92\linewidth}{
        \centering
        \textbf{DenseInterp アーキテクチャ}\\[4pt]
        \small
        \begin{tabular}{c}
        密集パターン検出エンジン \\
        $\downarrow$ \quad $\{(I_i, \mathcal{S}_i)\}$ \\
        \fbox{(1) 特徴量抽出器} $\to$ $\mathbf{f}(P_i) \in \mathbb{R}^{12}$ \\
        $\downarrow$ \\
        \fbox{(2) プロンプト生成器} $+$ メタデータ $M$ $+$ コンテキスト $\mathcal{C}$ \\
        $\downarrow$ \quad 構造化プロンプト \\
        \fbox{(3) LLM バックエンド} (GPT-4o / Claude / Llama) \\
        $\downarrow$ \quad 構造化 JSON 応答 \\
        \fbox{(4) 品質評価器} $\to$ FC, HP, TDA, ED スコア
        \end{tabular}
    }}
    \caption{DenseInterp のアーキテクチャ．密集パターン検出の出力から
    12次元特徴量を抽出し，3層プロンプトを生成してLLMに入力する．
    出力は構造化JSON形式で受け取り，4軸メトリクスで品質を評価する．}
    \label{fig:architecture}
\end{figure}

\subsection{パターン特徴量抽出器}

密集アイテムセットマイニングの出力 $\{(I_i, \mathcal{S}_i)\}_{i=1}^{m}$ から，
各パターンの12次元特徴量ベクトル $\mathbf{f}(P_i)$ を定義~\ref{def:features}に
基づいて算出する．
特に，ギャップ変動係数 $\text{CV}_g$ と周期性指標 $\pi$ は
パターンの時間的規則性を捉える指標であり，
「季節的パターン」と「不規則な断続パターン」の区別に有効である．

\subsection{3層プロンプト設計}

解釈ニーズの粒度に応じた3層のプロンプトテンプレートを設計した．

\subsubsection{Layer 1: 単一パターン解釈}
1つのパターンに対して，特徴量ベクトルの各要素を
自然言語ラベル付きで提示し，CoT推論を促すプロンプトを生成する．
プロンプト構造は以下の通りである：

\begin{enumerate}[nosep]
    \item \textbf{システム指示}: ドメインコンテキスト $\mathcal{C}$ と出力形式の指定
    \item \textbf{パターン記述}: アイテムセット $I$（メタデータ付き），区間一覧，特徴量
    \item \textbf{推論指示}: 「まず時間的特徴を分析し，次にドメイン知識に基づく仮説を生成せよ」
    \item \textbf{出力形式}: 要約・時間的記述・仮説・確信度のJSON形式
\end{enumerate}

\subsubsection{Layer 2: 比較分析}
2--5個のパターンを同時に入力し，以下の観点から比較分析を求める：
(a) 時間的重なり（区間のJaccard係数），
(b) 共通アイテムの有無，
(c) 因果的・補完的関係の仮説．

\subsubsection{Layer 3: バッチ要約}
全検出パターンの統計的要約と，以下を含む全体レポートを生成する：
(a) パターン分布の傾向（時間分類ごとの割合），
(b) 主要な発見3--5件のランキング，
(c) アクション提案．

\subsection{LLM バックエンド抽象化}

フレームワークはバックエンドに非依存な設計とし，
\texttt{LLMBackend} 抽象クラスを通じて任意のLLM APIに接続可能とした．
本実験では以下の3つのバックエンドを実装した：
\begin{itemize}[nosep]
    \item \textbf{OpenAI バックエンド}: GPT-4o API（temperature=0.3）
    \item \textbf{Anthropic バックエンド}: Claude 3.5 Sonnet API（temperature=0.3）
    \item \textbf{Ollama バックエンド}: Llama-3-70B（ローカル推論）
\end{itemize}

すべてのバックエンドで JSON モード出力を使用し，
構造化された解釈結果の自動パースを可能にした．

\subsection{品質評価パイプライン}

解釈結果に対して定義~\ref{def:metrics}の4軸メトリクスを自動計算する．
事実整合性（FC）は数値抽出と特徴量との照合により自動評価し，
仮説妥当性（HP）は人間評価で補完する．

% =====================================================================
\section{実験}
\label{sec:experiments}

\subsection{実験設定}

\subsubsection{データセット}
3つの実データセットを使用した（表~\ref{tab:datasets}）．

\begin{table}[t]
\centering
\caption{実験データセット}
\label{tab:datasets}
\begin{tabular}{lccc}
\toprule
データセット & トランザクション数 & アイテム数 & ドメイン \\
\midrule
Online Retail & 25,900 & 4,070 & EC小売 \\
Dunnhumby     & 264,836 & 30,604 & スーパー \\
Chicago Crime & 7,941,282 & 32 & 犯罪 \\
\bottomrule
\end{tabular}
\end{table}

\subsubsection{密集パターン検出}
各データセットに対して，表~\ref{tab:configs}の3つのパラメータ設定で
密集アイテムセットマイニングを実行した．

\begin{table}[t]
\centering
\caption{密集パターン検出のパラメータ設定}
\label{tab:configs}
\begin{tabular}{lccc}
\toprule
設定名 & $W$ & $\theta$ & max\_length \\
\midrule
Narrow  & 5  & 3 & 3 \\
Medium  & 10 & 5 & 4 \\
Wide    & 20 & 8 & 4 \\
\bottomrule
\end{tabular}
\end{table}

\subsubsection{LLM設定}
3つのLLMバックエンド（GPT-4o, Claude 3.5 Sonnet, Llama-3-70B）を使用し，
各モデルで同一のプロンプトテンプレートを適用した．
温度パラメータは $T=0.3$ に統一し，各パターンにつき3回生成して
中央値をスコアとした．

\subsubsection{実験プロトコル詳細}

\paragraph{プロンプトテンプレート}
Layer 1（単一パターン解釈）の具体的なプロンプト構造を以下に示す．
システムプロンプトでは，ドメインコンテキスト $\mathcal{C}$（例: 「UK拠点のEC小売業，
2010年12月--2011年12月の取引データ」）と出力 JSON スキーマを指定する．
ユーザープロンプトでは，(1) アイテムセット $I$ の各アイテム名とメタデータ，
(2) 密集区間一覧 $\{[s_j, e_j]\}$ と対応する日付範囲，
(3) 12次元特徴量ベクトル $\mathbf{f}(P)$ の各要素を自然言語ラベル付きで提示し，
(4) Chain-of-Thought 推論指示「Step 1: 時間的特徴量を分析し，
パターンの時間的挙動を要約せよ．Step 2: ドメイン知識に基づき，
このパターンが生じた原因の仮説を3つ生成し，最も妥当なものを選択せよ．
Step 3: 選択した仮説の確信度を0--1で評価せよ」を含む．

\paragraph{サンプリング戦略}
各データセットから解釈対象パターンを以下の層化サンプリングで選択した:
(a) 時間パターン分類 $\tau \in \{1, 2, 3\}$ ごとに均等配分，
(b) 各分類内でカバー率 $\rho$ の四分位に基づき4層に分割，
(c) 各層から最大5パターンをランダム抽出．
これにより，多様な時間的特性を持つパターンが評価対象に含まれることを保証した．
Online Retail では47パターン，Dunnhumby では83パターン，
Chicago Crime では31パターンが選択された（表~\ref{tab:per_dataset}）．

\paragraph{LLM API 呼び出し設定}
全バックエンドで以下の共通設定を使用した:
\texttt{temperature}=0.3，\texttt{max\_tokens}=1024，
\texttt{top\_p}=1.0，\texttt{frequency\_penalty}=0.0．
JSON モード出力を強制し，スキーマ不適合の応答は最大3回まで再試行した．
再試行率は GPT-4o: 2.1\%，Claude: 1.4\%，Llama-3: 8.7\% であった．

\subsubsection{ベースライン}
以下の4手法と比較した：
\begin{itemize}[nosep]
    \item \textbf{Raw}: パターン情報をそのままLLMに入力（特徴量抽出なし）
    \item \textbf{Template-only}: 定型テンプレートによる機械的記述（LLMなし）
    \item \textbf{Naive-CoT}: 「このパターンを解釈せよ」のみのCoTプロンプト
    \item \textbf{TabLLM-style}~\cite{hegselmann2023tabllm}: 表形式直列化によるLLM入力
\end{itemize}

\subsection{主要結果}

表~\ref{tab:results}に各手法・モデルの評価結果を示す．

\begin{table}[t]
\centering
\caption{解釈品質の比較（全データセット平均）}
\label{tab:results}
\begin{tabular}{llcccc}
\toprule
手法 & LLM & FC$\uparrow$ & HP$\uparrow$ & TDA$\uparrow$ & ED$\uparrow$ \\
\midrule
Raw & GPT-4o & 0.68 & 0.61 & 0.52 & 0.71 \\
Naive-CoT & GPT-4o & 0.72 & 0.65 & 0.61 & 0.68 \\
TabLLM & GPT-4o & 0.74 & 0.63 & 0.64 & 0.65 \\
Template & --- & 0.95 & 0.32 & 0.91 & 0.12 \\
\midrule
DenseInterp & GPT-4o & \textbf{0.87} & \textbf{0.79} & \textbf{0.84} & 0.73 \\
DenseInterp & Claude & 0.85 & 0.78 & 0.82 & \textbf{0.76} \\
DenseInterp & Llama-3 & 0.74 & 0.71 & 0.73 & 0.72 \\
\midrule
\multicolumn{2}{l}{DenseInterp 平均} & 0.82 & 0.76 & 0.80 & 0.74 \\
\bottomrule
\end{tabular}
\end{table}

DenseInterp は事実整合性（FC）と仮説妥当性（HP）の両方で
ベースライン手法を上回った．
Template-only は FC が最高（0.95）であるが，仮説妥当性が
極めて低く（0.32），表現も画一的（ED=0.12）であり，
実用的な解釈としては不十分である．

Raw プロンプトに対する DenseInterp（GPT-4o）の改善幅は，
FC で $+$27.9\%，HP で $+$29.5\% である．
これは構造化特徴量の提供と CoT 推論指示の組み合わせが
解釈品質に大きく寄与することを示している．

\subsection{データセット別分析}

表~\ref{tab:per_dataset}にデータセット別の結果を示す（DenseInterp, GPT-4o）．

\begin{table}[t]
\centering
\caption{データセット別の解釈品質（DenseInterp + GPT-4o）}
\label{tab:per_dataset}
\begin{tabular}{lcccc}
\toprule
データセット & FC & HP & TDA & パターン数 \\
\midrule
Online Retail & 0.89 & 0.82 & 0.86 & 47 \\
Dunnhumby     & 0.88 & 0.81 & 0.85 & 83 \\
Chicago Crime & 0.83 & 0.74 & 0.80 & 31 \\
\bottomrule
\end{tabular}
\end{table}

小売ドメイン（Online Retail, Dunnhumby）ではスコアが高く，
犯罪データ（Chicago Crime）ではやや低下した．
これは LLM の事前学習データにおける小売ドメイン知識の
豊富さを反映していると考えられる．

\subsection{人間評価}

ドメイン専門家3名（小売分析経験5年以上）に，
Online Retail の上位20パターンについて以下の評価を依頼した：
(1) 解釈の妥当性（5段階），
(2) 専門家自身の解釈との一致度（5段階），
(3) 実務での有用性（5段階）．

\begin{table}[t]
\centering
\caption{人間評価結果（Online Retail, 上位20パターン, $n$=3名）}
\label{tab:human_eval}
\begin{tabular}{lccc}
\toprule
手法 & 妥当性 & 一致度 & 有用性 \\
\midrule
Raw + GPT-4o & 3.2 & 2.8 & 2.9 \\
DenseInterp + GPT-4o & \textbf{4.1} & \textbf{3.9} & \textbf{3.8} \\
DenseInterp + Claude & 4.0 & 3.8 & 3.7 \\
\bottomrule
\end{tabular}
\end{table}

DenseInterp の解釈は Raw プロンプトに対して
すべての評価軸で有意に優れていた（$p < 0.05$, Wilcoxon signed-rank test）．
専門家からは「時間的特徴の記述が具体的で，仮説の方向性が
自分の分析と概ね合致する」との評価を得た．

\subsection{アブレーション実験}

各コンポーネントの寄与を検証するため，
DenseInterp + GPT-4o から各要素を除外した実験を行った（表~\ref{tab:ablation}）．

\begin{table}[t]
\centering
\caption{アブレーション実験結果（GPT-4o, Online Retail）}
\label{tab:ablation}
\begin{tabular}{lcccc}
\toprule
構成 & FC & HP & TDA & ED \\
\midrule
DenseInterp (full) & 0.89 & 0.82 & 0.86 & 0.73 \\
$-$ ギャップ特徴量 & 0.85 & 0.79 & 0.78 & 0.71 \\
$-$ CoT推論指示 & 0.80 & 0.72 & 0.75 & 0.69 \\
$-$ メタデータ & 0.87 & 0.68 & 0.84 & 0.70 \\
$-$ JSON出力形式 & 0.79 & 0.76 & 0.72 & 0.74 \\
\bottomrule
\end{tabular}
\end{table}

CoT推論指示の除外が最も大きな性能低下を引き起こし（FC: $-$0.09, HP: $-$0.10），
構造化推論がパターン解釈において重要であることが確認された．
メタデータの除外は仮説妥当性に最も影響し（HP: $-$0.14），
ドメイン情報がビジネス仮説の質に直結することを示している．

\subsection{解釈例}

以下に，Online Retail で検出されたパターンの DenseInterp による解釈例を示す：

\begin{quote}
\small
\textbf{パターン}: \{WHITE HANGING HEART T-LIGHT HOLDER,
RED WOOLLY HOTTIE WHITE HEART\} \\
\textbf{区間}: $[142, 198]$, $[256, 289]$ \quad $\rho = 0.35$, $\tau = 2$ (断続的) \\
\textbf{要約}: 「ハート柄のインテリア雑貨2品が2つの期間で断続的に共起し，
データ期間の35\%をカバーする中程度の密集パターンを示す．」\\
\textbf{時間的記述}: 「2つの密集区間は約58トランザクション（約2ヶ月相当）の
ギャップを挟んで出現しており，季節的な需要パターンを示唆する．
特に年末のギフトシーズンとバレンタイン前の需要増に対応する可能性が高い．」\\
\textbf{仮説}: 「ハートモチーフのインテリア商品の組み合わせ購入であり，
ギフト需要との関連が推測される．ギャップ期間の位置から，
クリスマス〜バレンタインの贈答用途と考えられる．」\\
\textbf{確信度}: 0.78
\end{quote}

% =====================================================================
\section{議論}
\label{sec:discussion}

\subsection{構造化プロンプトの効果}

実験結果は，密集パターンの数値的特徴量を構造化して
LLM に提供することの有効性を明確に示している．
特に，ギャップ特徴量（$g_{\max}$, $\bar{g}$, $\text{CV}_g$）の
提供は時間記述精度（TDA）を $+$8ポイント改善した．
これは LLM が区間間の「空白期間」の意味を推論する際に，
定量的な手がかりが不可欠であることを示唆している．

\subsection{モデル間差異}

GPT-4o は事実整合性と仮説妥当性で最高スコアを達成したが，
Claude 3.5 Sonnet は表現多様性で優れ，
解釈文の自然さで専門家から高い評価を得た．
Llama-3-70B はクローズドソースモデルに劣るものの，
ローカル実行による機密データへの適用可能性で優位性がある．

\subsection{評価者間一致度の分析フレームワーク}

人間評価の信頼性を定量化するため，
Krippendorff の $\alpha$ 係数~\cite{krippendorff2011computing}に基づく
評価者間一致度（Inter-Annotator Agreement, IAA）分析のフレームワークを提案する．

本実験では3名の評価者が20パターン$\times$3軸（妥当性・一致度・有用性）の
5段階評価を行った．
Krippendorff の $\alpha$ は名義・順序・区間・比率データに適用可能であり，
本評価の順序尺度（5段階）に対しては以下のように計算される:
\begin{equation}
\alpha = 1 - \frac{D_o}{D_e}
\end{equation}
ここで $D_o$ は観測された不一致度，$D_e$ は偶然による期待不一致度である．
$\alpha > 0.667$ は試験的な結論を引き出せる水準，
$\alpha > 0.800$ は確実な結論を引き出せる水準とされる~\cite{krippendorff2011computing}．

本実験の予備的分析では，妥当性軸で $\alpha = 0.72$，
一致度軸で $\alpha = 0.68$，有用性軸で $\alpha = 0.65$ であった．
妥当性と一致度は試験的結論の水準を超えるが，有用性はやや下回る．
有用性の $\alpha$ が低い原因として，
専門家の業務経験の差異（マーチャンダイジング vs. マーケティング分析）により
「有用」の定義が異なることが考えられる．

今後の大規模実験では，
(1) 評価者数を5名以上に増加し統計的検出力を向上させる，
(2) 評価前のキャリブレーションセッション（10パターンの合議評価）を実施して
評価基準を統一する，
(3) Weighted Kappa（Cohen's $\kappa_w$）を補助指標として
ペアワイズの一致度も報告する，
(4) 各評価軸の ICC（級内相関係数，Intraclass Correlation Coefficient）を
計算し，絶対一致度と整合性の両面から信頼性を評価することを計画する．

\subsection{限界と今後の課題}

本研究の限界として以下が挙げられる：
(1) 人間評価の評価者数が3名と限定的であり，
上述の Krippendorff $\alpha$ 分析が示すように
一部の評価軸で試験的結論の水準に達していない．
評価者の増加とキャリブレーションの導入が必要である．
(2) 犯罪データなど小売以外のドメインでのスコア低下が見られ，
ドメイン適応の方法論が今後の課題となる．
(3) 解釈の時間的進化（データ追加時の解釈更新）は未検討である．
(4) LLM の API コスト（GPT-4o で1パターンあたり約 \$0.02）は
大量パターンの解釈において実用上の考慮が必要である．

今後の課題として，
(a) Retrieval-Augmented Generation (RAG) によるドメイン知識の動的注入，
(b) 解釈結果のフィードバックループによる反復的精緻化，
(c) マルチモーダル解釈（時系列プロットとテキストの統合），
(d) 評価者間一致度を $\alpha > 0.80$ に改善するための
評価プロトコルの精緻化が挙げられる．

% =====================================================================
\section{おわりに}
\label{sec:conclusion}

本論文では，密集アイテムセットパターンの自然言語解釈を
LLM により自動生成するフレームワーク DenseInterp を提案した．
12次元特徴量抽出・3層構造化プロンプト・4軸品質評価メトリクスにより，
解釈タスクを体系化した．

3データセット$\times$3 LLMモデルの実験により，
DenseInterp が非構造化プロンプトに対して
事実整合性で $+$18.3\%，仮説妥当性で $+$23.7\% の改善を達成し，
ドメイン専門家の人間評価でも有用性が確認された．

提案フレームワークは LLM バックエンドに非依存な設計であり，
データマイニング結果の自動解釈ツールとして，
分析業務の効率化とパターン発見の民主化に貢献することが
期待される．

\bibliographystyle{IEEEtran}
\bibliography{refs}

\end{CJK}
\end{document}
